\documentclass[11pt,a4paper]{article}
\usepackage[utf8]{inputenc}
\usepackage[T1]{fontenc}
\usepackage{lmodern,amsmath,amssymb,amsthm,mathtools,array,booktabs,enumitem,microtype}
\usepackage{natbib}
\usepackage[margin=1in]{geometry}
\usepackage{xr-hyper}
\usepackage{hyperref}
\externaldocument[GE-]{01_law_of_graphic_equalisation_theory_referenced}
\externaldocument[PH-]{01a_graphic_equalisation_physics_grounding}
\externaldocument[NM-]{03_law_of_narrative_momentum_theory_referenced}
\newtheorem{definition}{Definition}[section]
\newtheorem{theorem}[definition]{Theorem}
\newtheorem{proposition}[definition]{Proposition}
\newtheorem{corollary}[definition]{Corollary}
\theoremstyle{remark}\newtheorem{remark}[definition]{Remark}
\title{The Law of Graphic Equalisation --- Born Recovery\\\large An Axiom-by-Axiom Note}
\author{Matthew Parslow\\Independent Researcher}
\date{\textbf{Draft \texorpdfstring{\textemdash}{---}\ draft-2026-09-02.1}\\[2pt]2 September 2026 \texorpdfstring{\textbullet}{\textperiodcentered}\ document revision 4}
\usepackage{fancyhdr}
\pagestyle{fancy}
\fancyhf{}
\fancyfoot[C]{\thepage}
\fancyfoot[R]{\scriptsize\textsc{draft-2026-09-02.1\,r4}}
\renewcommand{\headrulewidth}{0pt}
\fancypagestyle{plain}{\fancyhf{}\fancyfoot[C]{\thepage}%
  \fancyfoot[R]{\scriptsize\textsc{draft-2026-09-02.1\,r4}}\renewcommand{\headrulewidth}{0pt}}
\begin{document}
\maketitle
\begin{center}
\fbox{\parbox{0.92\textwidth}{\small
\textbf{DRAFT --- draft-2026-09-02.1 r4} (2 September 2026).
This document is under active revision and is circulated for comment, not as a
finished result. It is a companion to \emph{The Law of Graphic Equalisation ---
Theory} (referred to below as Graphic Equalisation) and to \emph{The Law of
Graphic Equalisation --- Physics Grounding} (the physics companion), and inherits
those papers' definitions, epistemic boundary, and level-of-claims; \emph{The Law of
Narrative Momentum} (the momentum paper) supplies the Narrative and the first law. It is an
\emph{attempt}, rated step by step, and not a peer-reviewed result. Sections
marked as falsifiability commitments are the intended points of attack.}}
\end{center}
\vspace{1em}

\begin{abstract}
This note checks, one axiom at a time, whether the change-substrate of Graphic
Equalisation supplies the operational axioms from which the informational
reconstructions of quantum theory derive the complex Hilbert space and the Born
rule. The papers named above carry the surrounding framework. The honest verb throughout is \emph{recover}, not \emph{derive}: the substrate
reproduces the known rule within its own structures, given its admitted
primitives and a contingent, self-located selection, rather than compelling it
from necessity. Each step is rated solid, earned, plausible or contingent, and the
residuals are named rather than hidden. Two findings carry the note. First, the substrate
instantiates the operational-probabilistic core --- convex states, affine
effects, a transformation monoid with reversibles, composition, normalised
probabilities --- from the classical total-probability backbone and not from any
quantum assumption. Second, two independent reconstructions have two
\emph{different} most-quantum axioms, and the substrate supplies each natively
from a different one of its own principles: Hardy's continuity of reversible
transformations from free continuation, and the purification of Chiribella,
D'Ariano and Perinotti from entanglement as a superframe dimension. Complex over
real is then not a gap to be closed by a forcing argument but the signature of a
contingent selection, self-located exactly as the number of spatial dimensions
is. The only claim made here is the \emph{arrangement} --- that the two crux axioms
above are already-present substrate structure; each component is individually known,
and the uniqueness is supplied by the cited theorems.
\end{abstract}

\section{The substrate, in one paragraph}\label{sec:substrate}
Seven notions carry the mapping below. Each is defined in one of the three papers
named above; this section fixes the references and adds nothing.
A \emph{frame} is a subgraph that runs (Graphic Equalisation's
Definition~\ref{GE-def:frame}), and frames compose by \emph{nesting}, a graph being a
subset of its container (its Proposition~\ref{GE-prop:no-direct-self-alter}).
A frame's state is its \emph{Narrative}, the standing structure it carries forward
(the momentum paper's Definition~\ref{NM-def:narrative}), and where a frame
instantiates a direction --- optional in general, present in this setting --- that
Narrative has a magnitude and a direction (the momentum paper's
Definition~\ref{NM-def:narrative-momentum}); the state proper is the \emph{internal}
Narrative, as against what another frame reads off it (its
Definition~\ref{NM-def:internal-external}).
A \emph{closure cycle} is the loop a realisation must close, the cycle being the unit
of inferred time (Graphic Equalisation's Proposition~\ref{GE-prop:cycle-unit}), and
\emph{phase} is position on that cycle rather than a signed time (its
Remark~\ref{GE-rem:phase-not-time}).
A \emph{realisation} is a collapse, a local write that picks one outcome from the
admissible set and carries a weight (Graphic Equalisation's
Definition~\ref{GE-def:collapse}), and
realisations are strictly ordered, the order carried by the rules rather than by a
clock (its Proposition~\ref{GE-prop:partial-order}).
\emph{Free continuation} is what a frame does when no rule fires, the
information-preserving default (the momentum paper's
Definition~\ref{NM-def:free-continuation}); \emph{free} means information-preserving
and not absent, since evaluating a rule is itself a write (Graphic Equalisation's
Proposition~\ref{GE-prop:evaluation-writes}), and all the mapping below uses is that
free continuation is \emph{injective}, hence reversible.
Nothing is privileged: there is no god's-eye frame, and therefore no privileged phase
origin and no privileged measurement basis (its Axiom~\ref{GE-ax:no-view}).
\emph{Entanglement} is a superframe dimension: a subframe's local state is mixed
because the purity is a dimension of the larger frame (the physics companion's
Proposition~\ref{PH-prop:mixed} and its Corollary~\ref{PH-cor:nonlocal}).

\section{The claim, honestly sized}\label{sec:claim}
The substrate \emph{recovers} the Born rule --- reproduces it within its own
structures, given its admitted primitives and a contingent, self-located
selection --- rather than \emph{deriving} it as a forced necessity.

\begin{remark}[Recovery is the honest verb, and the right one]\label{rem:recover}
Derivation per se would compel Born from primitives. The framework forces almost
nothing, and that is the point of a graph substrate: what survives every
attempted removal is a container relation across which rules run (Graphic
Equalisation's Remark~\ref{GE-rem:only-g}), and essentially everything else is
contingent rule-content, as with the number of spatial dimensions (the physics
companion's Section~\ref{PH-sec:dimensions}) or the gauge group (its
Remark~\ref{PH-rem:gauge-contingent}). A contingent law can only be recovered given its
rule-content, never derived as compelled. Some steps below are forced \emph{given
the premises} --- the squared modulus, from no-privilege together with the
additivity of decohered alternatives --- while the premises themselves are
admitted and contingent: complex, linear, self-similar.
\end{remark}

Every distinctively-quantum ingredient is supplied by a substrate principle; the
cited reconstruction theorems supply the uniqueness
\citep{gleason1957,hardy2001,masanes2011,chiribella2011}. The only real claim is
the \emph{arrangement} --- each component is individually known. That sizing is
not modesty about the result but a statement of what kind of result it is, and it
governs every rating below.

\section{Level 0: the substrate is an operational-probabilistic theory}\label{sec:level0}
Before any reconstruction axiom applies, the substrate must be an
operational-probabilistic theory (OPT): states in a convex set, effects,
transformations, and a composition rule. That framework is a \emph{precondition}
for either axiom set, so the mapping has two levels, and this is the lower one.

The key that makes it tractable is seeing what the OPT convex structure
\emph{is}. The convexity of the state space and the affinity of effects are the
\emph{classical total-probability backbone} --- total probability over ordered
branchings --- distinct from, and prior to, the quantum superposition linearity
that Section~\ref{sec:axioms} worries about. Every probabilistic theory has this
backbone. The substrate has it by construction, because realisations carry
weights, compose in order, and depend only on the present.

\begin{proposition}[The substrate instantiates the operational-probabilistic core]\label{prop:opt}
Convex states, affine effects, a transformation monoid with reversibles,
composition, and normalised probabilities are each supplied by the substrate's
weighted, ordered realisations together with the history-optional,
present-structure-acts first law (the momentum paper's Section~\ref{NM-sec:law};
Graphic Equalisation's Corollary~\ref{GE-cor:history}). The instantiation rests on the classical
total-probability backbone and on no quantum assumption.
\end{proposition}

Ingredient by ingredient, with the substrate structure that supplies each. Four
ratings are used here and in Section~\ref{sec:axioms}: \emph{solid}, the substrate
structure is the requirement's statement; \emph{earned}, derived here from a
substrate principle rather than read off one; \emph{plausible}, a commitment of the
right kind, with statement-match not shown; \emph{contingent}, rule-content, admitted
and self-located rather than supplied.

\begin{center}\small
\begin{tabular}{@{}>{\raggedright\arraybackslash}p{0.26\textwidth}
                >{\raggedright\arraybackslash}p{0.51\textwidth}
                >{\raggedright\arraybackslash}p{0.11\textwidth}@{}}
\toprule
\textbf{OPT requirement} & \textbf{Supplied by} & \textbf{Rating}\\
\midrule
Systems & Frame types: a subgraph together with an outcome set. & solid\\
States are well-defined equivalence classes & Statistics depend only on the current Narrative --- history is optional and present structure acts, so there is no preparation-memory and operational equivalence holds. & solid\\
Convex state space & Mixing two preparations by a weighted coin-realisation: realise a $\lambda$-weighted coin, prepare $P$ on one branch and $Q$ on the other. & solid\\
Affine effects & Total probability over the coin--prepare--measure branching, $e(\lambda P+(1-\lambda)Q)=\lambda e(P)+(1-\lambda)e(Q)$. & solid\\
Transformations (monoid, with reversibles) & Rules compose, in order; free continuation is the identity, information-preserving and hence reversible. & solid\\
Composition & Subframe nesting; a joint preparation lives over the product outcome set. Existence is Level 0; the complex-specific tomography is the contingent layer of Section~\ref{sec:complex}. & solid\\
Probabilities & Realisation weights in $[0,1]$, normalised over a complete measurement. & solid\\
\bottomrule
\end{tabular}
\end{center}

\begin{remark}[The affinity is not the superposition linearity]\label{rem:affine}
The affinity of effects is the law of total probability applied to ordered
realisations, and nothing about it is quantum. It is worth marking because the
two linearities are easy to run together, and the argument would be circular if
they were: Level 0 earns the classical backbone, and Section~\ref{sec:axioms}
must earn the superposition linearity separately, on other grounds.
\end{remark}

What remains at this level is hygiene rather than core: a finite tomographic
dimension, in the sense that a state is fixed by finitely many fiducial effects
--- plausible for finite frames by the aperture, unproven in general;
topological closure and compactness of the convex state space; the full
composition axioms, associativity of sequential and parallel composition and the
interchange law, which should follow from the ordering and nesting of realisations
but each need checking; and the no-restriction and effect-realisability
niceties. None of these is the question whether the state space convex-linearises
at all, which Level 0 answers; they are carried forward in
Section~\ref{sec:residuals}.

\section{The distinctively-quantum axioms, mapped}\label{sec:axioms}
Two independent reconstructions are used. \citet{hardy2001} gives five axioms, of
which classical probability theory satisfies the first four and fails the fifth.
\citet{chiribella2011} --- Chiribella, D'Ariano and Perinotti, CDP below --- give
six principles that single out quantum theory over the real, the quaternionic and
the classical. They have \emph{different} most-quantum axioms, which is what makes
the pair worth running rather than either alone.

\begin{proposition}[Each reconstruction's crux axiom is supplied natively, from a different substrate principle]\label{prop:crux}
Hardy's quantum-selecting axiom is the continuity of reversible transformations
between pure states; the substrate supplies it from free continuation, which is
information-preserving and hence reversible, and whose phase is a continuous
$U(1)$, so it moves continuously between pure states --- this note reads an
amplitude as $r\,e^{i\theta}$, modulus and cycle-phase. CDP's characteristically
quantum principle is purification; the substrate supplies it from entanglement as
a larger-frame dimension, a mixed subframe state being the restriction of a pure
superframe state. Neither is borrowed from the other, and neither was introduced
for this purpose.
\end{proposition}

That two separate reconstructions' crux axioms are both already-present substrate
structure, from two different principles, is the substantive positive finding of
this note. The remaining correspondences are weaker and are marked as such:

\begin{center}\small
\begin{tabular}{@{}>{\raggedright\arraybackslash}p{0.30\textwidth}
                >{\raggedright\arraybackslash}p{0.47\textwidth}
                >{\raggedright\arraybackslash}p{0.11\textwidth}@{}}
\toprule
\textbf{Axiom} & \textbf{Substrate structure} & \textbf{Rating}\\
\midrule
Continuity of reversible transformations (Hardy; classical theory fails exactly this) & Free continuation: information-preserving, hence reversible, with a continuous cycle-phase. & solid\\
Purification (CDP; fails classically) & Entanglement as a superframe dimension: a mixed subframe state is the restriction of a pure superframe state. & solid\\
Phase- and context-independence & No privileged phase origin and no privileged frame: probability is invariant under global phase and the same across measurement bases. & earned\\
Probability & Weighted realisations; frequencies over an ensemble of collapses from one possibility structure. & plausible\\
Simplicity and ideal compression & Minimum description length as a native commitment (Graphic Equalisation's Remark~\ref{GE-rem:mdl-ontic}) \citep{rissanen1978}. & plausible\\
Subspaces & Subframe-as-frame: restriction to a subset of distinctions behaves as a smaller frame. & plausible\\
Perfect distinguishability & Distinctions are first-class; a state that is not completely mixed carries one. & plausible\\
Pure conditioning & Collapse of a pure Narrative yields a definite realised outcome. & plausible\\
Local tomography (Hardy's composite axiom; CDP's local distinguishability) & Subframe composition over the product outcome set --- \emph{contingent}, and the subject of Section~\ref{sec:complex}. & contingent\\
\bottomrule
\end{tabular}
\end{center}

The plausible ratings are genre-matches rather than statement-matches: the
substrate has a commitment of the right kind, and that the specific axiom is the
one the substrate enforces is not shown. Marking them so is the point of the
table; a reader who upgrades one of them is disagreeing with a rating, which is
the intended form of the disagreement.

The squared modulus is selected among the natural candidates by phase-averaged
additivity, and made unique by Gleason, in the physics companion (its
Proposition~\ref{PH-prop:born} and the argument following it); this note uses the
result. What the framework supplies to Gleason \citep{gleason1957} are his premises
--- a vector-valued state, from the linearity step, and a context-independent
measure, from the absence of a privileged rule set --- and Gleason supplies what the
framework does not.

\begin{remark}[Two supporting identifications]\label{rem:identifications}
The \emph{inner product} is the cross-frame projection: one frame's Narrative
read in another's basis, so it is not imported structure but the framework's own
cross-frame reading. \emph{Unitarity} is that a free continuation destroys no
distinction --- an information-preserving write rather than the absence of one,
which is the right shape, since unitary evolution plainly does change the state
and is reversible for exactly that reason. Both are identifications rather than
derivations, and the rigorous step from ``uncoupled'' to full additivity and
norm-preservation is the standing residual of Section~\ref{sec:residuals}.
\end{remark}

\section{Complex over real: contingent, not forced}\label{sec:complex}
Local tomography is the axiom that selects complex Hilbert space over real, and it
is where a forcing argument would have to sit. Writing $K$ for Hardy's fiducial
degrees of freedom including normalisation, for two two-state systems:

\begin{center}\small
\begin{tabular}{@{}lcccl@{}}
\toprule
\textbf{Theory} & $K$ \textbf{single} & $K$ \textbf{composite} & $K_A\!\cdot\!K_B$ & \textbf{Local tomography}\\
\midrule
Classical & 2 & 4 & 4 & holds (but fails continuity)\\
Real (rebit) & 3 & 10 & 9 & \textbf{fails} ($10>9$)\\
Complex (qubit) & 4 & 16 & 16 & \textbf{holds}\\
\bottomrule
\end{tabular}
\end{center}

A closure-phase frame is \emph{consistent with} complex amplitudes at one system
--- the Bloch sphere, $K=4$, rather than the real disk's $K=3$ --- though a
phase-less frame is equally constructible and the substrate requires neither.
Read that way, a superframe, being a frame, is a complex closure-phase Narrative
over the product outcome set, giving $K_{AB}=16$: local-tomographic, complex.

\begin{proposition}[Complex over real is a self-located selection, not a forced consequence]\label{prop:contingent}
Complex-over-real is contingent rule-content,
as with the number of spatial dimensions. Local tomography is therefore not
a gap to be closed by a forcing argument; it is the signature of a contingent
selection. A real-amplitude composition, $K_{AB}=10$, is another admissible
configuration. Our universe is the complex one, self-located (the physics
companion's Proposition~\ref{PH-prop:self-location}).
\end{proposition}

\begin{remark}[Why forcing is the wrong question]\label{rem:misframing}
Whether the substrate \emph{forces} the complex composite is not the open question,
and asking it would be an error of the framework's own kind rather than merely an
unproven step. Seeking to force complex over real would contradict the framework's stance that
it never forces contingent physics. The reading that survives is narrower and
does real work: the admissible complex configuration yields $K_{AB}=16$ and the
admissible real one yields $10$; a real composite would be a superframe that is
not a closure-phase frame like its subframes, so real amplitudes cost
self-similarity. That gives a direction and a reason without giving a compulsion,
and it is the same status every other contingent constant in this programme has.
\end{remark}

Given the contingent choice, the Born rule and the Hilbert structure follow via
the mapping and the cited reconstructions --- which is how the framework treats
all physics: shapes follow from contingent rule-content, the rule-content itself
never forced. Of the post-2011 reconstructions, \citet{masanes2011} obtain
finite-dimensional quantum theory from five physical requirements on states and
transformations, and theirs is the formulation this note would most like tested
against, since those requirements are stated on single systems plus a composition
rule --- exactly the seam where the substrate's answer turns contingent rather than
forced. Independent support for reading the selection as contingent comes from
\citet{hoehn2017} and \citet{hoehnwever2017}, whose rules on an observer's
acquisition of information yield the correlation structure of qubits and rebits
alike and so do not by themselves select complex over real. The rest of the
reconstruction landscape is the physics companion's
(its Section~\ref{PH-sec:prior-art}).

\section{What this reconciles}\label{sec:reconcile}
Special and general relativity reconcile on this substrate as one change-budget
spent two ways, with the invariance of $c$ a definitional consequence rather than a
postulate, and quantum collapse and relativity are joined at $c$ as the in-frame
event and the cross-frame map of one change-substrate (the physics companion's
Proposition~\ref{PH-prop:sr-gr}, its Corollary~\ref{PH-cor:reconcile}, and its
Section~\ref{PH-sec:collapse}). With Born \emph{recovered} on the same primitives,
Born, SR and GR stand reconciled on one substrate --- the reconciliation is that
companion's, argued in its Section~\ref{PH-sec:born-sr-gr}, and it is a
reconciliation of the \emph{conflict} at the level of shape, not a quantum theory
of gravity: no field is quantised, no metric derived, no dimensionless constant
fixed.

\section{Residuals}\label{sec:residuals}
What is open, unaddressed, and named:

\begin{enumerate}[leftmargin=*]
\item \textbf{Level 0 hygiene.} Whether ``present structure acts'' genuinely
delivers operational equivalence --- states as equivalence classes of preparations
with identical statistics --- and whether the full OPT composition axioms
(associativity of sequential and parallel composition, the interchange law) hold
under the ordering and nesting of realisations. With them, a finite tomographic
dimension, topological closure and compactness of the state space, and
no-restriction and effect-realisability. These are hygiene rather than core, and
plausible for finite frames; none is proven here.
\item \textbf{The linearity step.} The rigorous passage from ``uncoupled free
continuation'' to full additivity and norm-preservation. The mechanism motivates
linearity; it does not prove it, and the linearity of the imposed supergraph map
is not handed over for free. This is the one posit the recovery rests on.
\item \textbf{The squaring argument.} Whether the phase-average selection is tight
enough on its own, or whether the full Gleason content is needed to reach
uniqueness. As the physics companion states it (its Proposition~\ref{PH-prop:born}
and the argument following it) the phase average selects rather than proves, and
Gleason is cited for the rest.
\item \textbf{Composition.} Under the natural reading the composite is
local-tomographic and complex. Whether that reading is forced, permitted, or
wrong is not settled here. The formulation of \citet{masanes2011} is the one this
seam should be tested against, its requirements being stated on single systems
plus a composition rule; that is exactly where the substrate's answer turns
contingent rather than forced, and where the mapping is least confident that it is
doing honest work.
\item \textbf{The plausible ratings.} Probability, simplicity and compression,
subspaces, distinguishability and pure conditioning are genre-matches to native
commitments, not statement-matches. Each would need tightening before the cited
theorems could be applied without qualification.
\end{enumerate}

\section{Falsifiability commitments}\label{sec:falsifiability}
Each commitment below is marked \textsc{[shape]} or \textsc{[detail]}. A
\textsc{[detail]} failing is a revision: the account survives with that mechanism
replaced. A \textsc{[shape]} failing costs \emph{range} rather than a mechanism,
because what fails is something the substrate is claimed to require or to permit
--- so the account does not hold where it claimed to, which is the more serious of
the two and is still a boundary rather than an annihilation.

\begin{enumerate}[leftmargin=*]
\item \textsc{[shape]}~\textbf{The operational-probabilistic instantiation.} Fails
if the substrate does not instantiate an OPT --- if convex states, affine effects,
a transformation monoid with reversibles, composition and normalised probabilities
cannot all be had from weighted ordered realisations and the history-optional first
law. Everything below is conditional on this, so it fails with it.
\item \textsc{[detail]}~\textbf{The Level-0 hygiene conditions.} Fails if the
conditions of Section~\ref{sec:residuals} --- a finite tomographic dimension,
topological closure and compactness of the state space, the full composition axioms
--- cannot be met for the frames this note is about. Each is a mechanism the note
could lose and replace without losing the instantiation above.
\item \textsc{[detail]}~\textbf{Linearity from free continuation.} Fails if the
step from uncoupled independent continuation to full additivity and
norm-preservation cannot be made rigorous --- in which case the vector structure is
a posit rather than a recovery, and Gleason's first premise is imported rather than
supplied.
\item \textsc{[shape]}~\textbf{The axiom mapping.} Fails if it does not go through
for a named axiom: in particular if free continuation is not Hardy's continuous
reversible transformation between pure states, or if entanglement as a superframe
dimension is not CDP's purification. Failure at either of those two is worse than
failure elsewhere, since they are the axioms that separate quantum from classical
and are what Proposition~\ref{prop:crux} claims.
\item \textsc{[detail]}~\textbf{The squared modulus.} Fails if, given complex
linearly superposing amplitudes and no privileged phase, the probability is not the
squared modulus; and the phase-averaged selection argument of the physics companion
(its Proposition~\ref{PH-prop:born} and the argument following it) fails separately
if another natural candidate satisfies the phase-averaged additivity, leaving the
whole weight on Gleason.
\item \textsc{[detail]}~\textbf{The composition seam.} Fails if the substrate's
composition rule, tested against requirements stated on single systems plus a
composition rule \citep{masanes2011}, yields a composite that is neither the complex
admissible configuration ($K_{AB}=16$) nor the real one ($K_{AB}=10$) --- which
would make the contingency verdict of Proposition~\ref{prop:contingent} a verdict
about the wrong set.
\item \textsc{[shape]}~\textbf{Complex over real is contingent.} Fails if the
framework must \emph{force} complex over real rather than admit both and
self-locate --- that is, if a substrate that admits a real-amplitude composition is
thereby shown to be the wrong substrate, rather than one whose selection among
admissible configurations is a matter of where the observer is.
\end{enumerate}

\noindent\emph{Scope note (not a falsifier).} This note is shape-level and not
quantitative: it fixes no dimensionless constant, derives no field equation, and
offers no theory of quantum gravity. Nor is it a derivation: the claim is the
\emph{arrangement}, and no step above is offered as compelling the Born rule from
primitives.

\begin{thebibliography}{99}

\bibitem[Chiribella et~al.(2011)]{chiribella2011}
Giulio Chiribella, Giacomo Mauro D'Ariano, and Paolo Perinotti.
\newblock Informational derivation of quantum theory.
\newblock \emph{Physical Review A}, 84(1):012311, 2011.

\bibitem[Gleason(1957)]{gleason1957}
Andrew M. Gleason.
\newblock Measures on the closed subspaces of a Hilbert space.
\newblock \emph{Journal of Mathematics and Mechanics}, 6:885--893, 1957.

\bibitem[Hardy(2001)]{hardy2001}
Lucien Hardy.
\newblock Quantum theory from five reasonable axioms.
\newblock arXiv:quant-ph/0101012, 2001.

\bibitem[H\"ohn(2017)]{hoehn2017}
Philipp A. H\"ohn.
\newblock Toolbox for reconstructing quantum theory from rules on information acquisition.
\newblock \emph{Quantum}, 1:38, 2017.

\bibitem[H\"ohn and Wever(2017)]{hoehnwever2017}
Philipp A. H\"ohn and Christopher S.~P. Wever.
\newblock Quantum theory from questions.
\newblock \emph{Physical Review A}, 95(1):012102, 2017.

\bibitem[Masanes and M\"uller(2011)]{masanes2011}
Llu\'is Masanes and Markus P. M\"uller.
\newblock A derivation of quantum theory from physical requirements.
\newblock \emph{New Journal of Physics}, 13(6):063001, 2011.

\bibitem[Rissanen(1978)]{rissanen1978}
Jorma Rissanen.
\newblock Modeling by shortest data description.
\newblock \emph{Automatica}, 14(5):465--471, 1978.

\end{thebibliography}
\end{document}
